\begin{table}[H]
\centering\centering
\caption{Race gap in veteran-contract pay \label{tab:pay-gap-veteran}}
\centering
\resizebox{\ifdim\width>\linewidth\linewidth\else\width\fi}{!}{
\begin{tabular}[t]{lccccccccc}
\toprule
  & (1) & (2) & (3) & (4) & (5) & (6) & (7) & (8) & (9)\\
\midrule
P(Black) & -0.023 & -0.015 & -0.026 & -0.035 & -0.004 & -0.039 & -0.039 & -0.001 & 0.098\\
 & (0.070) & (0.058) & (0.049) & (0.049) & (0.051) & (0.045) & (0.046) & (0.055) & (0.073)\\
P(other race) & -0.072 & -0.131 & -0.087 & -0.073 & -0.038 & -0.096 & -0.092 & 0.012 & -0.019\\
 & (0.123) & (0.101) & (0.081) & (0.081) & (0.085) & (0.076) & (0.079) & (0.090) & (0.131)\\
\midrule
Observations & 5062 & 5062 & 5062 & 5062 & 5062 & 5062 & 5062 & 5062 & 5062\\
R$^2$ & 0.236 & 0.439 & 0.673 & 0.690 & 0.717 & 0.755 & 0.784 & 0.686 & 0.139\\
Within R$^2$ & 0.150 & 0.376 & 0.636 & 0.655 & 0.685 & 0.727 & 0.732 & 0.651 & 0.042\\
Position $\times$ year FE & Yes & Yes & Yes & Yes & Yes & Yes & Yes & Yes & Yes\\
Career stage (A) & No & Yes & Yes & Yes & Yes & Yes & Yes & Yes & No\\
Prior-season production (B) & No & No & Yes & Yes & Yes & Yes & Yes & Yes & No\\
Career production (C) & No & No & No & Yes & Yes & Yes & Yes & Yes & No\\
Pre-NFL signals (D) & No & No & No & No & Yes & Yes & Yes & Yes & No\\
Usage (E) & No & No & No & No & No & Yes & Yes & No & No\\
Signing team $\times$ year FE & No & No & No & No & No & No & Yes & No & No\\
Draft slot, round, extension dummy & Bucket (prior) & Ext. + bucket & Ext. + bucket & Ext. + bucket & Yes & Yes & Yes & No & No\\
Race-prior covariate FE & Yes & Yes & Yes & Yes & Yes & Yes & Yes & No draft & No draft\\
Mean of outcome & 1.154 & 1.154 & 1.154 & 1.154 & 1.154 & 1.154 & 1.154 & 1.154 & 1.154\\
Implied Black gap (\%) & -2.3 & -1.5 & -2.5 & -3.4 & -0.4 & -3.8 & -3.9 & -0.1 & 10.2\\
SE of P(Black) clustered by player and signing team & (0.061) & (0.050) & (0.044) & (0.042) & (0.042) & (0.038) & (0.038) & (0.046) & (0.060)\\
Residual SD of P(Black) given column's regressors & 0.275 & 0.273 & 0.269 & 0.265 & 0.243 & 0.239 & 0.228 & 0.245 & 0.278\\
\bottomrule
\multicolumn{10}{l}{\rule{0pt}{1em}* p $<$ 0.1, ** p $<$ 0.05, *** p $<$ 0.01}\\
\end{tabular}}
\end{table}

{\footnotesize
\noindent\textit{Notes:} This table includes the estimation results of equation (1). The outcome is log APY (\$ millions); with OTC position $\times$ year-signed fixed effects this equals log APY over the league salary cap of the signing year, without OTC's rounding of the cap share. The sample is freely bargained veteran contracts (unrestricted free agents, including those who re-sign with their own team, and extensions) signed 2014-2026 by players with a race prediction; tags and tenders, whose pay is set by a CBA formula, are in Table \ref{tab:pay-gap-terms}. 0 contracts of players with no race prediction and 0 with missing APY are excluded. Every column includes OTC position (18 market positions) $\times$ year-signed fixed effects; column (7) adds signing franchise $\times$ year-signed fixed effects. Column (8) is column (5) without draft slot, draft-round dummies and the extension dummy. Every column also includes fixed effects for the covariates of the race prediction's prior (position at NFL entry, entry era, draft-round bucket, college type and whether a home county is known), as regression calibration requires. In column (8), P(Black) and P(other race) come from the draft-free prediction (a prior without the draft-round bucket), and the draft-round bucket is dropped from the prior fixed effects, so no regressor records the draft. Because the prior's covariates include the draft-round bucket, columns (1)-(4) already condition on draft round and column (1) is not a raw gap. Column (9) is the raw benchmark: P(Black) and P(other race) from the draft-free prediction, with its prior's covariates (no draft bucket) and position $\times$ year-signed fixed effects only; the difference between columns (9) and (1) is the part of the raw gap that conditioning on draft round removes. Career stage (A): experience-bin dummies (0, 1, 2, 3, 4-6, 7+ seasons), age at signing and its square, a left-censored-experience indicator and an extension dummy (reference: unrestricted free agents). Whether a team extends a player is itself a team decision. Prior-season production (B): season-before-signing box-score and PFR measures with position-group-specific slopes (passing for QBs, rushing and receiving for RBs, receiving for WRs and TEs, tackles, sacks, QB hits, tackles for loss and pressures for DL and LB, tackles, interceptions, passes defended and coverage allowed for DBs, kicking and punting), games played and injury weeks interacted with position group, and a no-prior-season indicator. Career production (C): career games and injury weeks interacted with position group, and career sums of the block B production measures with the same position-group slopes. Pre-NFL signals (D): log draft pick, an undrafted indicator, draft-round dummies, 247 rating and stars, combine measures (40, vertical, bench, broad jump, cone, shuttle, height, weight) and athletic scores (a composite and size, speed, explosion and agility subscores built from the combine measures) interacted with position group, and the final college team's Power-conference, SRS and HBCU indicators. Usage (E): prior-season offensive, defensive and special-teams snaps, prior-season games started, career depth-chart starts and career snaps, interacted with position group. Snaps and starts are chosen by coaches, so the usage controls may absorb discrimination in playing time. Controls with incomplete coverage are set to zero with a missing indicator (interacted like the control). Games played and the box-score volume counts (attempts, targets, tackles) also depend on playing time, and draft slot and the extension dummy record earlier team decisions, so columns (3)-(5) estimate a gap conditional on earlier allocation decisions, which discrimination could also affect. The identifying assumption is that, conditional on the controls and fixed effects, race is uncorrelated with unobserved productivity. Omitted quality would bias the estimate: there are no PFF-type grades, and offensive linemen have no production measure, so their quality rests on games, injury weeks and pre-NFL signals. Regression calibration identifies the Black-white gap only if, in addition to names and hometown being unrelated to pay given race and the controls, (i) P(Black) is calibrated given every control in the column, not only given the prior's covariates, and (ii) the gap does not vary with the controls (otherwise the coefficient is a variance-weighted average of gaps). Controls that predict P(Black) within the prior's cells do not by themselves violate (i): under (i) they predict P(Black) exactly when they predict race. They do reduce the variation in P(Black) that identifies the coefficient (overlap), and if (i) fails for them the coefficient moves as they are added for reasons unrelated to pay. Table \ref{tab:pay-gap-veteran} reports the identifying variation by column and checks (i) against documented race. The row 'Residual SD of P(Black) given column's regressors' is the standard deviation of P(Black) after partialling out each column's fixed effects, prior-covariate dummies and controls, the variation that identifies the coefficient on P(Black); it falls from 0.275 in column (1) to 0.243 in column (5) (the column-(5) regressors explain 62.0\% of the variance of P(Black)). This measures overlap and precision, not calibration: under condition (i) the controls predict P(Black) exactly when they predict race. Condition (i) is checked against documented race on the 1,435 contracts of 762 players whose race a public source states (28.3\% of the sample; 1,092 documented non-Hispanic Black alone, the event P(Black) models, and 343 of any other documented race, documented Hispanic and multiracial Black players among them as non-target; conflicting sources excluded). A logit of the documented label on logit P(Black), the prior-covariate dummies and position fixed effects (player-clustered) has slope 0.647 (SE 0.064) (one under calibrated relative odds). Adding the controls' P(Black) index ($X'\hat\gamma$: the column-(5) controls weighted by their coefficients in a regression of P(Black) on the column-(5) regressors, a function of observed characteristics and P(Black) only, never of pay) gives 0.733 (SE 1.108), $p$ = 0.508; under (i) it is zero. This is one restriction in the one direction along which the controls predict P(Black): a non-rejection does not establish calibration given every control, and a rejection shows that the documented label departs from P(Black) along that direction. Documentation is positive-only and requires a Wikipedia article: documented players are mostly Black and famous (mean P(Black) 0.693 against 0.578 for the undocumented), selection on race shifts only the logit intercept, and selection on fame within race is not controlled, so the check describes documented players and makes no claim about calibration in the population. P(Black) is the predicted probability of being non-Hispanic Black alone; Black Hispanic and multiracial Black players enter P(other race), and Table \ref{tab:pay-gap-race-measures} also reports estimates with P(Black) + P(multiracial) as the Black regressor. The predicted Black share is not calibrated in levels: it lies below the TIDES African-American player share through 2016 and above the self-identified Black-alone share (below the Black-or-multiracial share) from 2019 (Table \ref{tab:race-prediction-tides}); Table \ref{tab:pay-gap-race-measures} reports estimates with the TIDES-raked posterior. Standard errors treat the predicted probabilities as known; the prior is estimated by EM on the whole player population of the database, so the omitted first-stage variance is likely small. Column (5) is the main specification. The implied gap is $100(e^{\hat\beta_1}-1)$. The 95\% confidence interval for P(Black) in column (5) is [-0.104, 0.096] log points; the residual standard deviation of P(Black) given the column-(5) regressors is 0.243 (row 'Residual SD of P(Black)'), so the estimates are imprecise. Standard errors are clustered at the player level. The row 'SE of P(Black) clustered by player and signing team' reports the standard error of the race coefficient with two-way clustering by player and signing franchise (teams set pay). P(Black) and P(other race) enter together, so the coefficient on P(Black) compares a Black (predicted: non-Hispanic Black alone) with a white player. Race is predicted, not observed: each person's probability of being non-Hispanic Black alone combines first name, surname and hometown county (BIFSG) with an NFL-specific prior estimated by EM on predetermined characteristics (players: position at entry, rookie era, draft round, college type, county availability; staff: role, unit and era at first appearance); documented race statements are not used (notes/race-prediction-design.md). Black Hispanic and multiracial Black persons are in the other-race probability, so P(Black) is narrower than the Black-alone-or-in-combination definition of the hand codes and of TIDES. Regressions use the probabilities (regression calibration) and control for the prior's covariates; the coefficient on P(Black) is the Black-white gap if the probabilities are calibrated given the regression's controls and names and hometown are unrelated to the outcome given race and the controls. The validation exhibit compares the predictions with published TIDES shares.
\par}

